\documentclass[11pt,a4paper]{article}
\usepackage{amsmath,amssymb,geometry,graphicx,booktabs,hyperref,microtype,enumitem}
\geometry{margin=2.5cm}
\hypersetup{colorlinks=true,linkcolor=blue,citecolor=blue,urlcolor=blue}

\title{\textbf{Paper 51: Necessity as Method}\\
\large Bottom-Up Derivation of Primitives in Adaptive Memory Research\\
\normalsize Papers~0--50}
\author{VCSM Project}
\date{2026-03-10}

\begin{document}
\maketitle

\begin{abstract}
We describe the methodological framework that produced Papers~1--50
and argue that it is structurally identical to the framework it produced.
The core pattern: begin with an engineering constraint rather than a hypothesis;
let necessity reduce the mechanism to its primitive; gate theoretical consolidation
on experimental survival; stress-test every layer adversarially.
We show that this process is a literal instance of VCSM's three-timescale learning
rule (session $\to$ paper $\to$ framework document)
with viability gating (only findings that survive adversarial stress consolidate),
mandatory turnover (every paper terminates when its question closes),
and copy-forward inheritance (the framework document seeds each new paper's
starting state via memory files).
\textbf{The meta-claim}: the research process exhibits zone identity encoding ---
a topographic, non-parity understanding of the theory ---
because it ran the full VCML formation process.
Reading the framework document without running the process
produces parity encoding at best (Paper~50).
The project is a self-demonstrating proof.
\end{abstract}

\section{Introduction}

Fifty papers and approximately 200{,}000 experiments later, the VCSM framework
has a precise, falsified-and-refined account of how viability-constrained coupled
map lattices build spatial memory through mandatory individual turnover.
The question that motivated this paper is simpler: \emph{how was it built?}

Not ``what is the theory?''\ but ``what kind of research process produces theory
of this kind?'' --- theory that is load-bearing, minimal, and adversarially robust.

The answer, it turns out, is structural rather than procedural.
The research process is not merely analogous to the VCSM framework.
It is an instance of it.

This paper makes that claim precise, provides evidence from the project history,
and draws out a consequence that is not obvious: \emph{zone identity encoding in
the theory is not transmissible} (Paper~50) --- it requires independent formation,
not relay from the framework document.
The self-demonstrating structure is therefore not merely an aesthetic observation.
It is a constraint on how the theory can be understood, extended, or replicated.

\section{Six Methodological Principles}

\subsection{Start with a constraint, not a hypothesis}

A hypothesis predicts what is true.
A constraint states what the system \emph{must be able to do.}
These are different starting points with different failure modes.

VCSM began with: \emph{cells need to die, or the system crystallises.}
This is not a hypothesis to be confirmed --- it is a structural requirement
derived from the viability constraint itself.
If the field only accumulates, it locks.
If it only decays, nothing consolidates.
Mandatory turnover is the \emph{only} mechanism that simultaneously satisfies
both the adaptivity requirement (old structure can be overwritten) and
the persistence requirement (new structure consolidates before the cell dies).

The constraint forces the mechanism.
Starting from ``I think turnover is useful'' would have been epistemically weaker:
it invites confirmation bias, permits optional death, and misses the
engineering necessity that makes death load-bearing rather than incidental.

The same pattern recurs throughout the project.
Paper~12 asked not ``does the copy-forward loop improve performance?''
but ``what prevents window collapse in extreme conditions?''
The copy-forward loop was the only mechanism that could --- and it was found
by following the constraint, not by predicting it.

\subsection{Let necessity do the reduction}

Primitives are not chosen top-down.
They are what is left after necessity prunes everything contingent.

VCSM's four primitives (viability constraint, copy-forward loop,
VCSM learning rule, zone-differentiated wave environment) were not
postulated in Paper~1.
They accumulated one by one as each ablation study removed something
that turned out to be dispensable:
frequency encoding without structured waves (dispensable),
GRU hidden state (dispensable at HS=1, Paper~80),
global synchronisation (dispensable),
quiescence gate for formation under continuous waves
(dispensable --- Paper~43 found formation is environmentally driven, not rule-specific).

What could not be dispensed with is what is real.
This is the constraint-first analogue of Occam's razor, but operationalised
experimentally rather than aesthetically.

\subsection{Bottom-up experimental ladder}

Papers do not follow a predetermined plan.
Each paper closes one question and opens the next.
The sequence is a viability-gated exploration:
failed hypotheses are the signal, not noise.

Eleven hypotheses did not survive:
\begin{itemize}
  \item \textbf{Paper~9}: A single governing ratio $\Xi$ collapses the adaptive window. NULL.
  \item \textbf{Paper~14}: Two dimensionless ratios $(R_A, R_B)$ fully parameterise sg4. FAIL.
  \item \textbf{Paper~15}: sg4 follows a power law in amplitude. REVISED (saturation curve;
        power law was a mixed-phase artifact).
  \item \textbf{Paper~35/V73}: Dynamic turnover always outperforms static reservoirs. REVERSED
        (static wins below the critical zone width; dynamic advantage only holds when
        $\delta \geq 10$~columns and wave density is appropriate).
  \item \textbf{Paper~41}: Task switching below $T_\mathrm{crit} \approx 800$ steps causes
        net memory loss. NULL (every switching period amplifies structure).
  \item \textbf{Paper~43}: Viability gating is necessary for formation.
        REVERSED (formation is environmentally driven; gate is an adversarial filter).
  \item \textbf{Paper~49 (original)}: mean\_relay (copy-forward seeding from L1) outperforms
        ctrl. REVERSED with full metric suite (ctrl wins; mean\_relay birth injection washed out).
  \item \textbf{Paper~79/V79}: Cascade input attenuates structure. REVERSED
        ($2.57\times$ more sg4 under cascade; hierarchy amplifies, not abstracts).
  \item \textbf{Paper~49/V90}: Mobius sign structure (period-2 pattern across generations).
        NULL (sign consistency $= 0.510$, chance).
  \item \textbf{Papers~35--39 (interpretation)}: Geographic relay preserves zone structure.
        QUALIFIED (relay gains in sg4 are real, but structure is parity encoding, na$<1$,
        not identity encoding; qualified by Papers~49 and~50).
  \item \textbf{Paper~49 (relay)}: Any relay coupling can recover identity encoding in L2.
        NULL (Paper~50: only independent WaveEnvStd achieves na$>1$).
\end{itemize}

Each of these failures was informative.
The V73 reversal revealed the critical zone width threshold that anchors the
entire C$_2$ coordinate system.
The Paper~43 reversal separated formation (environmentally driven) from
maintenance (rule-governed), which is now a first-order distinction in the framework.
The Paper~49 relay reversal, via the na\_ratio metric introduced in Paper~47,
revealed the parity/identity distinction that Paper~50 closes.

The ladder is not linear. It is viability-gated: each paper is a cell that
lives long enough to produce its finding, then terminates.

\subsection{Adversarial stress-testing at every layer}

Every confirmed finding was subjected to adversarial challenge before
it consolidated into the framework document.
This occurred at three levels:

\textbf{Mechanism ablation.}
Paper~43 systematically ablated each component of the VCSM rule
(viability gate, copy-forward loop, zone-differentiated input, VCSM update)
and measured what broke.
Formation under zone-structured waves survived ablation of the gate.
Persistence (maintenance of structure through turnover) required the
copy-forward loop --- without it, structure degraded $1.63\times$ under
maintenance protocol.

\textbf{Environmental adversarial.}
Papers~46 and~47 directly tested the quiescence gate under adversarial
perturbation (non-zone-correlated, high-amplitude noise).
This revealed the gate reversal at $t^*$: gate is not a formation tool,
it is an adversarial filter, and the optimal strategy reverses at the
commitment epoch.

\textbf{Metric stress-testing.}
Paper~48 found that sg4 (the primary metric through Paper~47) gives
the wrong answer on the $\delta$ sweep: sg4 increases as $\delta$ shrinks
(a counting artifact), while decode\_norm and SNR correctly track causal direction.
This prompted the full metric suite mandate: every experiment from Paper~48
onward reports sg4, $\sigma_w$, SNR, na\_ratio, decode\_norm, mm\_lda, and coll/site.
The GPT-2 experiment series validated the causal geometry conjecture
against a rotation attack: principal component distances track ablation influence
only when accounting for sign flips at the feature-construction/causal-integration transition.

\subsection{Derive metrics from theoretical primitives}

A convenience metric measures what is easy to measure.
A theory-derived metric measures what the theory says is the operative quantity.

sg4 (zone-mean L2 distance) was the primary metric for Papers~1--47.
It is easy to compute and reflects genuine zone differentiation.
But Paper~48 showed it is systematically blind to the \emph{direction} of
causal influence, and Paper~49/50 showed it is blind to the \emph{type}
of zone encoding (identity vs.\ parity).

The metric evolution follows the theory:
\begin{enumerate}[label=\arabic*.]
  \item sg4 --- zone-mean signal (sufficient for Papers~1--40)
  \item $\sigma_w$, SNR --- noise floor and signal-to-noise (required by C$_2$ gate analysis)
  \item na\_ratio (nonadj/adj) --- zone topology, identity vs.\ parity (required by relay analysis)
  \item decode\_norm --- per-cell location accuracy (required by adversarial gate analysis)
  \item mm\_lda, coll/site --- mid-memory structure and turnover rate
\end{enumerate}

Each metric was added when the theory outgrew the previous one.
The metric is not a measurement of the theory; it is a question the theory asks.

\subsection{Maintain three timescales in the research process itself}

The VCSM learning rule operates on three timescales:
\begin{enumerate}
  \item \textbf{Calm (fast):} $\mathrm{baseline}_h[i] \mathrel{+}= \beta(\mathrm{hid}[i] - \mathrm{baseline}_h[i])$
  \item \textbf{Perturb (medium):} $\mathrm{mid\_mem}[i] \mathrel{+}= \alpha_\mathrm{mid}(\mathrm{hid}[i] - \mathrm{baseline}_h[i])$
  \item \textbf{Consolidate (slow, gated):} $\mathrm{fieldM}[i] \mathrel{+}= \alpha_\mathrm{field}(\mathrm{mid\_mem}[i] - \mathrm{fieldM}[i])$
\end{enumerate}

The research process operates on three timescales with identical structure:
\begin{enumerate}
  \item \textbf{Session (fast):} Experiments run; immediate results observed;
        session context discarded on completion.
        Analogous to baseline\_h tracking: high-frequency, lossy, non-consolidated.
  \item \textbf{Paper (medium):} Multiple sessions compose a paper.
        Findings are written to \texttt{MEMORY.md} and topic files.
        Analogous to mid\_mem: accumulated contrast signal, pending consolidation,
        subject to decay if the paper's question is not closed.
  \item \textbf{Framework document (slow, gated):} Findings that survive adversarial
        review consolidate into \texttt{VCML\_THEORY\_COMPLETE.tex}.
        Analogous to fieldM: stable, long-timescale, seeds each new paper's starting state
        via the copy-forward mechanism (new papers read the theory document and memory files).
\end{enumerate}

The quiescence gate in the research process is \emph{adversarial challenge}:
a finding does not enter the framework document until it has been stress-tested
at a minimum by a replication with changed parameters, an ablation,
and if possible a direct contradiction attempt.

\section{The Three-Timescale Mapping}

Table~\ref{tab:mapping} makes the structural homology explicit.

\begin{table}[h]
\centering
\caption{VCSM components and their research-process analogs.}
\label{tab:mapping}
\begin{tabular}{p{4.5cm}p{4.5cm}p{4.5cm}}
\toprule
\textbf{VCSM component} & \textbf{Research-process analog} & \textbf{Concrete instantiation} \\
\midrule
Cell & Experimental paper & Papers~1--50 \\
Viability constraint & Hypothesis must produce finding & Failed hypotheses terminate the paper \\
Mandatory turnover & Every paper terminates & Questions closed, not left open indefinitely \\
Calm state (baseline\_h) & Session working memory & Context window; discarded on session end \\
Perturbation (mid\_mem) & Paper-level accumulation & MEMORY.md, topic files; pending consolidation \\
Quiescence gate (SS$\geq$10) & Adversarial stress test & Must survive ablation / metric challenge \\
Consolidation (fieldM) & Framework document & VCML\_THEORY\_COMPLETE.tex \\
Copy-forward loop & Theory document seeds new paper & New paper reads framework + memory files \\
Wave environment & Experimental constraint space & Parameter sweeps, adversarial conditions \\
Zone-differentiated input & Structured research constraints & C$_2$ coordinates, metric suite mandate \\
Zone identity encoding & Topographic theory structure & na\_ratio$>1$ in theory space \\
\bottomrule
\end{tabular}
\end{table}

The analogy is not metaphorical at the level of ``research is like learning.''
It is structural: the same three-timescale architecture, the same gating criterion,
the same copy-forward inheritance, and --- critically --- the same result:
\emph{zone identity encoding that is not transmissible via relay}.

\section{The Relay Problem as Conclusion}

Paper~50 found that zone identity encoding requires the full VCML formation
process running independently in the receiving module.
No relay --- amplitude copy, graded amplitude, birth seeding, or any
combination --- achieves na\_ratio$>1$ in L2.
The identity-specific information was generated by L1's formation process and
lives in L1's fieldM.
It was never in the outgoing wave signal.

The identical structure holds for the VCSM research framework.

A reader of \texttt{VCML\_THEORY\_COMPLETE.tex} in isolation receives the
framework document the way L2 receives L1's wave amplitude array in geo\_copy relay:
an instantaneous snapshot of the consolidated signal, without the temporal
integration history.
That reader can achieve parity encoding --- a valid, useful understanding
of the even/odd structure of the theory (what is suppressed, what is excited,
which hypotheses died) --- but not identity encoding (the topographic structure
of \emph{why} each element is where it is, which constraints forced which primitives,
which failures were informative in which direction).

The identity encoding lives in the formation history: the 50-paper viability-gated
ladder, the 11 null results and their specific failure modes, the metric evolution
from sg4 to na\_ratio, the adversarial challenges and what they broke.
That history cannot be transmitted by reading the output.
It requires independent formation --- running the process, not receiving its fieldM.

This is not a limitation specific to VCSM.
It is a structural property of any theory whose content is predominantly
\emph{what was eliminated and why}.
A theory that accumulated by constraint propagation has most of its information
in the pruning history.
The surviving document shows the topology; it does not show the gradient.

\section{Discussion}

\subsection{Implications for replication}

The non-transmissibility of zone identity encoding implies that
``replication'' in the standard sense (re-running Paper~50's experiments
with the same code and parameters) produces parity encoding, not identity encoding.
A replicator achieves na\_ratio$=1.219$ for ctrl\_std but may not understand
\emph{why} ctrl\_std uniquely achieves identity, because that understanding
requires Paper~49's relay failure and Paper~47's na\_ratio derivation
and Paper~43's formation/maintenance separation.

True replication of the framework requires running the full experimental ladder.
This is not a methodological deficiency --- it is a structural property of
constraint-first research.
A framework built by constraint propagation can only be fully understood by
running the constraint propagation yourself.

\subsection{Implications for AI research methodology}

The six principles described in Section~2 are not specific to VCML.
They describe the conditions under which any empirical research process
can achieve theory of the constraint-propagated, adversarially-robust kind:

\begin{enumerate}[label=\arabic*.]
  \item Start from engineering necessity, not from hypothesis
  \item Reduce primitives by ablation, not by postulation
  \item Gate theoretical consolidation on adversarial survival
  \item Derive metrics from theoretical primitives, not from convenience
  \item Maintain three timescales: session, paper, framework
  \item Accept null results and reversals as the primary signal
\end{enumerate}

The dominant alternative --- hypothesis-driven research with predetermined
metrics and a predetermined ladder --- produces theories that are
``relay encoding'': structurally coherent but topographically impoverished,
because the identity information was never in the hypothesis to begin with.
The hypothesis encodes only what the researcher already understood.
The constraint encodes what the system requires.

\subsection{The Gödel connection}

A top-down critique of bottom-up methodology is self-refuting by construction.

To argue that the constraint-first, bottom-up method is wrong, one must argue
from a position of knowing in advance what the primitives are --- which is
exactly the top-down stance the methodology predicts will miss the primitive.
A critic who argues that VCSM's mandatory turnover is unnecessary must either
(a) have derived the turnover necessity from first principles (in which case
they are doing constraint propagation) or (b) be asserting from a theoretical
prior (in which case they are arguing from hypothesis, and the methodology
predicts they will miss the mechanism).

This is not a rhetorical trick.
It is structural.
The Gödel parallel: self-referential statements are gated out by the
quiescence condition (a system chasing its own tail never settles,
never writes to fieldM; see Paper~46's Gödel interpretation of the
identity-fluid property).
A self-refuting critique has the same topology --- it must use the method
to reject the method, and in doing so confirms it.

The methodology paper cannot be cleanly attacked from outside the methodology.
This is not a bug. It is the correct structure for a theory whose
content is primarily what was eliminated and why.

\subsection{The three-document architecture}

This paper is part of a three-document structure:
\begin{itemize}
  \item \textbf{VCML\_THEORY\_COMPLETE.tex} --- the \emph{what}: consolidated findings,
        mechanisms, equations, experimental evidence.
        Operates on the slow (fieldM) timescale.
  \item \textbf{Paper~51 (this paper)} --- the \emph{how}: methodology, structure of
        the research process, why the process produces what it produces.
        Operates on the medium (paper) timescale.
  \item \textbf{Papers~1--50} --- the \emph{why}: the experimental ladder, hypothesis deaths,
        metric derivations, adversarial challenges.
        Operates on the fast (session) timescale, accumulated into MEMORY.md.
\end{itemize}

Three documents. Three timescales. Same architecture as the framework they describe.

\subsection{The centerpiece example: parity/identity invisible for 15 papers}

The strongest evidence that the methodology is working is the relay finding.

Papers~35--49 reported geographic relay gains of $3$--$5\times$ in sg4.
The parity/identity distinction --- that relay produces zone \emph{parity} encoding
(na\_ratio$<1$, adjacent zones more separated than non-adjacent) rather than zone
\emph{identity} encoding (na\_ratio$>1$) --- was invisible for approximately 15 papers
because sg4 is positive for both.

The distinction became visible in one step when na\_ratio was derived from the
theoretical primitive (zone topography requires non-adjacent zones to be more
distinct than adjacent zones) rather than from measurement convenience.
Paper~47 derived the coordinate. Paper~49 applied it. Paper~50 closed the result.

This is the metric principle working exactly as described:
a convenience metric (sg4) asks whether zones are separated.
A theory-derived metric (na\_ratio) asks whether the \emph{correct} zones are separated.
The relay architecture looked successful for 15 papers because the question was wrong.

The correct question was invisible until the theory outgrew the metric.
The metric outgrew the theory only because the theory was built by constraint
propagation (what the relay \emph{has to} produce for downstream identity encoding)
rather than by hypothesis (relay probably preserves zone structure).

\subsection{The self-demonstrating structure}

The claim that the methodology is self-demonstrating is not merely aesthetic.
It is falsifiable.

If the methodology does not exhibit the properties it describes, the claim fails.
The properties it describes are:
(1) mandatory turnover (every paper terminates);
(2) viability gating (only adversarially-tested findings consolidate);
(3) three-timescale organisation;
(4) copy-forward inheritance;
(5) zone identity encoding (topographic, not just parity, structure in the theory).

Evidence for (1)--(4) is given in Sections~2--3.
Evidence for (5) is the structure of the framework document itself:
\texttt{VCML\_THEORY\_COMPLETE.tex} is organised not as a list of confirmed findings
(parity) but as a causal gradient from constraint through mechanism through
experimental evidence to prediction (identity).
The na\_ratio of the framework document, if one could measure it, is $>1$.

\section{Conclusion}

The VCSM research process is an instance of the VCSM framework.
It runs on three timescales (session, paper, framework document),
gates consolidation on adversarial survival (the quiescence condition),
enforces mandatory turnover (every paper terminates when its question closes),
and uses copy-forward inheritance (new papers read the accumulated theory).

The result is zone identity encoding --- a topographic, gradient-structured
theory of viability-constrained spatial memory --- that cannot be transmitted
by reading the output.

The framework is self-demonstrating not because we designed it that way,
but because the constraint that produced the framework also constrained
the process that produced it.
Engineering necessity, followed consistently, produces this structure
whether the substrate is a 3200-cell VCML lattice or a 50-paper research program.

The conclusion of Paper~50 is also the conclusion of this paper:
identity encoding requires the full formation process run independently.
You cannot learn this framework by reading it.
You have to build it.

\end{document}
